\section*{Asesmen}
\addcontentsline{toc}{section}{Asesmen}

Asesmen Bab XII mencakup Praktikum 3 (12\%) dan UTS Praktik (14\%).
Fokusnya adalah tabel ACTION/GOTO SLR dan LALR, serta kemampuan
menguji dan men-\textit{debug} konflik.

\begin{table}[H]
\centering
\caption{Instrumen asesmen Bab XII}
\label{tab:ch12-asesmen}
\begin{tabularx}{\textwidth}{@{}lXl@{}}
\toprule
\textbf{Instrumen} & \textbf{Indikator} & \textbf{Bobot} \\
\midrule
Praktikum 3 & Tabel ACTION/GOTO SLR dan LALR, log debug konflik, dan
uji string & 12\% (bagian) \\
UTS Praktik & Tabel SLR/LALR dan penyelesaian konflik shift-reduce
atau reduce-reduce & 14\% (bagian) \\
Kuis Konsep & Item LR, closure, goto, dan jenis konflik & formatif \\
\bottomrule
\end{tabularx}
\end{table}

Kriteria kelulusan: mahasiswa dapat membangun tabel ACTION/GOTO SLR
dan LALR, membaca laporan konflik, men-\textit{debug} konflik sampai
selesai, serta memilih strategi parsing yang sesuai.
